% ===== PRÉAMBULE CANONIQUE — à coller VERBATIM en tête de CHAQUE .tex =====
\documentclass[11pt,a4paper]{article}
\usepackage[utf8]{inputenc}\usepackage[T1]{fontenc}\usepackage{lmodern}
\usepackage[french]{babel}
\usepackage{amsmath,amssymb,mathtools}\usepackage{icomma}
\usepackage[version=4]{mhchem}          % \ce{...} : équations & formules
\usepackage{siunitx}\sisetup{locale=FR} % \qty{}{}, \si{}, \ang{}
\usepackage{chemfig}                    % \chemfig{...} : structures développées
\usepackage{xcolor}\usepackage{colortbl}
\usepackage{tikz}\usepackage{pgfplots}\pgfplotsset{compat=1.18}
\usepackage[most]{tcolorbox}\usepackage{enumitem}\usepackage{array,booktabs}
\usepackage[a4paper,margin=2.2cm]{geometry}
\usetikzlibrary{arrows.meta,calc,angles,quotes,positioning,patterns,backgrounds,shapes.geometric}
\definecolor{primary}{RGB}{222,49,99}\definecolor{primarydark}{RGB}{150,22,55}
\definecolor{defcol}{RGB}{0,114,178}\definecolor{methcol}{RGB}{52,92,148}
\definecolor{excol}{RGB}{0,135,90}\definecolor{attncol}{RGB}{200,40,55}
\tcbset{boxbase/.style={enhanced,breakable,boxrule=0.9pt,arc=2.5pt,
  left=3mm,right=3mm,top=2.3mm,bottom=2.3mm,fonttitle=\bfseries,coltitle=white}}
% --- boîtes TUTORIEL (noms EXACTS, ne pas renommer) ---
\newtcolorbox{cle}[1][]{boxbase,colback=defcol!6,colframe=defcol,
  title={\textsc{À retenir}\ifx\relax#1\relax\relax\else\textmd{ : \itshape #1}\fi}}
\newtcolorbox{methode}[1][]{boxbase,colback=methcol!7,colframe=methcol,sharp corners,
  title={\textsc{Méthode}\ifx\relax#1\relax\relax\else\textmd{ --- \itshape #1}\fi}}
\newtcolorbox{exemplebox}[1][]{boxbase,colback=excol!5,colframe=excol,
  title={\textsc{Exemple}\ifx\relax#1\relax\relax\else\textmd{ : \itshape #1}\fi}}
\newtcolorbox{piege}[1][]{boxbase,colback=attncol!6,colframe=attncol,
  title={\textsc{Piège}\ifx\relax#1\relax\relax\else\textmd{ : \itshape #1}\fi}}
\newtcolorbox{remarque}[1][]{boxbase,colback=black!4,colframe=black!45,
  title={\textsc{Remarque}\ifx\relax#1\relax\relax\else\textmd{ : \itshape #1}\fi}}
% --- boîtes EXERCICE / CORRIGÉ (noms EXACTS, auto-comptées) ---
\newtcolorbox[auto counter]{exo}[1][]{boxbase,colback=primary!4,colframe=primary,
  title={\textsc{Exercice}~\thetcbcounter\ifx\relax#1\relax\relax\else\textmd{ --- \itshape #1}\fi}}
\newtcolorbox[auto counter]{corrige}[1][]{boxbase,colback=excol!4,colframe=excol,
  title={\textsc{Corrigé}~\thetcbcounter\ifx\relax#1\relax\relax\else\textmd{ --- \itshape #1}\fi}}
\newcommand{\R}{\mathbb{R}}\newcommand{\N}{\mathbb{N}}\newcommand{\Z}{\mathbb{Z}}\newcommand{\ds}{\displaystyle}
% (un \usepackage{fancyhdr}+en-tête est possible ; il est ignoré côté web)
% =========================================================================
\begin{document}

\begin{exo}[reconnaître la loi]
Pour chaque situation, indique la grandeur maintenue constante et écris la relation à utiliser (sans
calcul).
\begin{enumerate}[label=\alph*)]
  \item On comprime un gaz à température constante.
  \item On chauffe un gaz dans un cylindre à piston libre (pression atmosphérique).
  \item On chauffe un gaz enfermé dans une bouteille rigide.
  \item On ajoute du gaz dans un ballon souple, à température et pression ambiantes constantes.
\end{enumerate}
\end{exo}

\begin{exo}[raisonnement proportionnel sans calcul détaillé]
Déduis l'effet demandé par proportionnalité (directe ou inverse).
\begin{enumerate}[label=\alph*)]
  \item À $V$ constant, la température absolue passe de \qty{1100}{\kelvin} à \qty{550}{\kelvin} : que
        devient la pression ?
  \item À $p$ constante, on double la température absolue : que devient le volume ?
  \item À $T$ constante, on divise la pression par deux : que devient le volume ?
  \item À $p$ constante, on double le volume : que devient la température absolue ?
\end{enumerate}
\end{exo}

\begin{exo}[loi de Charles avec conversions]
Un gaz est maintenu à pression constante. Convertis en kelvins puis calcule.
\begin{enumerate}[label=\alph*)]
  \item $V_1=\qty{2.0}{\litre}$ à \qty{27}{\degreeCelsius}, chauffé à \qty{87}{\degreeCelsius} : $V_2$ ?
  \item $V_1=\qty{4.0}{\litre}$ à \qty{127}{\degreeCelsius}, refroidi à \qty{27}{\degreeCelsius} : $V_2$ ?
  \item Le volume passe de \qty{5.0}{\litre} à \qty{6.0}{\litre} en partant de \qty{250}{\kelvin} :
        quelle est la température finale $T_2$ ?
\end{enumerate}
\end{exo}

\begin{exo}[isochore et Avogadro]
Applique la relation adaptée (attention aux kelvins pour les températures).
\begin{enumerate}[label=\alph*)]
  \item ($V$ constant) $p_1=\qty{1.0}{atm}$ à \qty{250}{\kelvin}, chauffé à \qty{500}{\kelvin} : $p_2$ ?
  \item ($V$ constant) La pression passe de \qty{2.0}{atm} à \qty{3.0}{atm} en partant de
        \qty{300}{\kelvin} : quelle est $T_2$ ?
  \item ($p,T$ constants) $V_1=\qty{6.0}{\litre}$ pour $n_1=\qty{0.25}{\mole}$ ; on porte la quantité à
        $n_2=\qty{0.40}{\mole}$ : $V_2$ ?
\end{enumerate}
\end{exo}

\end{document}
